% ---------------------------------------------------------------------------
% Section draft: the reward's structure (comfort mean + ventilation-bound bill).
% Standalone: guarded macros, so it compiles inside paper/main.tex and inside
% paper/sections_preview.tex.
% ---------------------------------------------------------------------------
\providecommand{\pending}[1]{\textcolor{red}{[pending: #1]}}
\providecommand{\e}[1]{\times10^{-#1}}

\section{Reward anatomy}
\label{sec:anatomy}

\paragraph{Comfort is a mean over the people present.} For zone $v$
with live band $[T^{\text{lo}}_v,T^{\text{hi}}_v]$ and shortfall $s_{v,t}$ (linear to
$5$\,K, quadratic beyond), the productivity term of the reward is
\begin{equation}
\frac{P-P^{\min}}{P^{\max}-P^{\min}}-1
\;=\;-\,\frac{\sum_v o_{v,t}\,s_{v,t}}{\sum_v o_{v,t}},
\label{eq:mean}
\end{equation}
the occupancy-weighted \emph{mean} shortfall over the occupants present
\citep{goldfeder2026smart}. Because it is a mean, the headcount cancels: a kelvin of shortfall
costs the same with two people in the building as with sixty, so the first arrivals of the
morning cost as much as a full building, and an empty building costs nothing. Occupancy is
exogenous: every zone holds the same headcount, people arrive between 06:00 and 12:00 and leave
between 13:00 and 19:00, and a controller that only reacts to arrivals pays at every morning
mode change. % src: docs/sim_analysis.md L192-198, L341-343
The base controller is more reactive than its parameters suggest: its preheat ramp sits inside a
flag that is true only in the hour before the mode change, so its ramps never exceeded one hour.
Moving the ramp outside the flag at a 2\,h horizon cuts comfort regret in the first occupied hour
fivefold and improves the score by $1.0\,\%$; releasing a building that is still empty at noon
adds $0.77\,\%$ more, independently. Longer horizons keep improving that hour but not the score:
readiness saturates once every zone is in band, while the gas keeps growing. % src: results/{refwu_dev30_shipped,ctrw_all,ctrw_d3,ctrw_d4}

\paragraph{The bill is envelope heat times a ventilation multiplier.} On the anchor
(\texttt{building\_train\_1/fairbanks\_1}, 2\,016 steps, mean outside $-26.3\,^\circ$C) the same
controller loses $93.12$ units, almost all of them to energy and carbon; comfort is $0.6\,\%$.
% src: results/houssem_anchor_ctrlw; docs/opt_gap.md §3
Instrumenting the trajectory splits the bill: about three quarters is reheat that replaces the
heat the envelope loses, just under a quarter is the outside-air share of that same reheat, and
fans, the boiler's standing loss and the pump make up the last few percent
(Appendix~\ref{app:energy}). % src: results/ef_shipped/building_train_1__fairbanks_1.npz (2016 steps)
The outside-air share is fixed by the action space: recirculation is capped at $r\le0.7$, so at
least $30\,\%$ of the supply air comes from outside, and the coil must lift it from
$T_{\text{mix}}=0.7\,T_{\text{ret}}+0.3\,T_{\text{out}}$ to the discharge cap
$T_{\text{dis}}=308.15$\,K, while the zone receives only the drop from $T_{\text{dis}}$ to
$T_{\text{zone}}$. With $T_{\text{ret}}\approx T_{\text{zone}}$ this gives
\begin{equation}
\frac{\text{gas billed}}{\text{net zone heat}}\;=\;
1+\frac{0.3\,\big(T_{\text{zone}}-T_{\text{out}}\big)}{T_{\text{dis}}-T_{\text{zone}}},
\label{eq:mult}
\end{equation}
so at the anchor's ambient the multiplier is $2.00$ for a zone held at $294$\,K and $1.18$ for one
at $270$\,K (Figure~\ref{fig:mult}a): \emph{heat delivered into a colder building is cheaper}. A
controller cannot change the $0.3$, but it can choose the temperature at which it delivers the
morning recovery. Beyond that, the air side leaves little to gain: over the episode CtrlW pays
only about $1\,\%$ more gas per delivered watt than a coil at its cap would.
% src: docs/control_gap.md §0.2, §2 (eff.py: gas/Q_net exact identity); docs/opt_gap.md §2-3

\paragraph{What the anatomy rules out.} The exchange rate between the two terms rules out
trading comfort in occupied hours for energy. On the anchor $624$ of the $2{,}016$ steps have
someone in the building, and a uniform $1$\,K shortfall in those steps costs
$0.6\times(1/5)\times624\approx75$ units, while lowering the whole building by one kelvin saves
about $0.57$\,MWh of gas, or $1.9$ units: a ratio of about $39{:}1$. What remains is the time
nobody is there. The night costs nothing in comfort, yet CtrlW, like the family it comes from,
heats it.
% src: docs/control_gap.md §3; docs/opt_gap.md §6; docs/energy_flow.md "By time bucket"
